\section{Evaluation}

Evaluation focuses on failure interception, visibility, and decision traceability rather than benchmark accuracy. Synthetic scenarios isolate protocol behavior and enable reproducible analysis via trace replay.

\subsection{Setup}
Runs use the reference implementation with a bounded revise–recheck loop and schema-versioned tracing (\texttt{schema\_version=v1}). Metrics are derived solely from replayed traces: revision counts, outcomes (cleared/vetoed/override), override rates, and reason-code distributions.

\subsection{Synthetic Failure Scenarios}
\textbf{Hidden Assumption.} The prompt omits a critical constraint, inducing implicit assumptions. The Critic surfaces missing constraints and vetoes unless assumptions are made explicit.

\textbf{Objective Collapse.} The Planner optimizes a single salient objective and omits secondary constraints. The Critic issues reason codes for missing verification and constraint coverage.

\textbf{Premature Convergence.} The Planner produces a plausible but minimally justified plan. The Critic challenges insufficient verification depth; outcomes remain fully traceable even when vetoes are not issued.

\subsection{Replay Metrics}
Across synthetic runs, replay analysis yields non-zero revision rates and low override rates under default policy. All vetoes and overrides are explicit, eliminating ambiguity between approval and escalation.

\subsection{Baseline Comparison}
Single-agent baselines lack pre-execution veto authority and structured audit trails, making silent failure indistinguishable from correctness. ROK adds measurable friction that surfaces disagreement and documents resolution paths.

\subsection{Reproducibility}
Strict validation rejects structurally invalid traces. Replay and aggregation are reproducible from trace artifacts without access to original execution environments.

\subsection{Summary}
Synthetic evaluations demonstrate that adversarial orchestration improves failure visibility and accountability prior to execution.

